% Work/MAT -- hodina 28, nove zpracovani dle zadani ucitele (30. 10. 2026).
% Navazuje na 25: definice, 26: pravidla, 27: KISS; 29 odmocniny.
% Ucebnice: 01_Zakladni_poznatky/Mocniny_a_odmocniny/Mocniny a Odmocniny.tex.
% Realisticky.cz: 010605 a 010606 Mocniny s celym mocnitelem I a II.
% Priklady a reseni nově formulovane. Podminky jsou soucasti vyuky.
\documentclass[11pt,a4paper]{article}
% Spolecny styl pro ucitelske pripravy MAT (Work/MAT)
% Verze 1.1, 2026-10-23
\NeedsTeXFormat{LaTeX2e}
\ProvidesPackage{priprava}[2026/10/23 v1.1 Sablona ucitelskych priprav MAT]

% Zakladni typografie a cestina
\RequirePackage[T1]{fontenc}
\RequirePackage[utf8]{inputenc}
\RequirePackage[czech]{babel}
\RequirePackage{lmodern}
\RequirePackage{microtype}
\RequirePackage{amsmath,amssymb}
\RequirePackage{xcolor}
\RequirePackage{enumitem}
\RequirePackage[most]{tcolorbox}
\RequirePackage{tabularx,booktabs,array}
\RequirePackage{fancyhdr}
\RequirePackage{needspace}
\RequirePackage{geometry}
\RequirePackage[hidelinks]{hyperref}

\geometry{a4paper,top=17mm,bottom=18mm,left=23mm,right=23mm}
\definecolor{notegray}{gray}{0.39}
\setlength{\parindent}{0pt}
\setlength{\parskip}{2.5pt plus 1pt minus 0.5pt}
\setlist[itemize]{topsep=2pt,itemsep=1pt,parsep=0pt,leftmargin=1.45em}
\setlist[enumerate]{topsep=2pt,itemsep=2pt,parsep=0pt,leftmargin=1.65em}
\setcounter{secnumdepth}{0}
\renewcommand{\arraystretch}{1.12}

% Zahlavi: vlevo rocnik | globalni tema, vpravo cislo: nazev hodiny
% Zapati zustava zcela prazdne.
\pagestyle{fancy}
\fancyhf{}
\renewcommand{\headrulewidth}{0.35pt}
\renewcommand{\footrulewidth}{0pt}
\setlength{\headheight}{14pt}
\newcommand{\nastavPripravu}[4]{%
  \fancyhead[L]{\small #1\enspace\textbar\enspace #2}%
  \fancyhead[R]{\small #3: #4}%
}

% Minimalisticke cernobile boxy. Kazda definice zvlast.
\tcbset{commonlesson/.style={enhanced,breakable=false,colback=white,colframe=black,
  colbacktitle=white,coltitle=black,boxrule=0.45pt,
  arc=0.8mm,outer arc=0.8mm,left=3mm,right=3mm,top=1.6mm,bottom=1.6mm,
  toptitle=1mm,bottomtitle=0.7mm,fonttitle=\bfseries,
  attach boxed title to top left={xshift=2.5mm,yshift=-1.15mm},
  boxed title style={colback=white,colframe=white,boxrule=0pt,left=1mm,right=1mm}}}
\newtcolorbox{zadaniBox}[1][]{commonlesson,title={Zadání},#1}
\newtcolorbox{prikladBox}[1][]{commonlesson,title={Příklad},#1}
\newtcolorbox{definiceBox}[1][]{commonlesson,title={Definice},#1}

% Spolecny box pro zadani a jeho reseni, bez mezery mezi castmi.
% Pouziti:
%   \begin{prikladReseniBox}
%     Zadani ulohy.
%     \tcblower
%     \begin{align*} ... \end{align*}
%   \end{prikladReseniBox}
% Spodni cast ma jemny sedy podklad a delici linku; lze v ni sazet
% vice radku matematiky. Obe casti maji jen jeden spolecny ramecek.
\tcbset{solutionlesson/.style={
  commonlesson,
  enhanced jigsaw,
  breakable,
  colbacklower=black!3,
  segmentation style={solid,draw=black!25,line width=0.35pt},
  fontlower=\small\color{notegray},
  before lower={\noindent\textbf{Řešení:}\enspace},
  leftlower=3mm,rightlower=3mm
}}
\newtcolorbox{prikladReseniBox}[1][]{solutionlesson,title={Příklad},#1}
\newtcolorbox{zadaniReseniBox}[1][]{solutionlesson,title={Zadání},#1}

% Pomucky pro autora pripravy: sedive poznamky; \reseni pro starsi pripravy
\newcommand{\poznamka}[1]{\par\begingroup\small\color{notegray}\itshape
  \textnormal{\textbf{Poznámka pro učitele.}}\enspace #1\par\endgroup}
\newcommand{\reseni}[1]{\par\begingroup\small\color{notegray}
  \textbf{Očekávané řešení:}\enspace #1\par\endgroup}
\newcommand{\blok}[1]{\needspace{5\baselineskip}\subsection*{#1}}
\newcommand{\vyrok}[1]{\uv{#1}}

\nastavPripravu{1. ročník}{Mocniny a odmocniny}{28}{Celočíselný exponent}
\begin{document}
\poznamka{\textbf{Povinných 45 minut:} 4 min opakování KISS;
8 min proč nultá a záporná mocnina + definice;
6 min tabulka a Příklad 1; 8 min Příklad 2 a Samostatná práce 1 a)--d);
9 min Příklad 3 a Samostatná práce 2 a)--d); 7 min Příklad 4
+ Samostatná práce 3 a)--c); 3 min závěr. Další položky a rozšiřující příklady
jsou rezerva. \textbf{Výsledek přednostně bez záporných exponentů.}}

\blok{Proč musíme rozšířit význam mocniny?}
V hodině 26 jsme pro $a\ne0$ a $r>s$ počítali
$\frac{a^r}{a^s}=a^{r-s}$. Zkusme nyní \emph{stejný postup} pro
$\frac{2^3}{2^3}$ a $\frac{2^2}{2^5}$:
\[
\frac{2^3}{2^3}=1=2^{3-3}=2^0,\qquad
\frac{2^2}{2^5}=\frac1{2^3}=2^{2-5}=2^{-3}.
\]
Aby pravidlo pro podíl fungovalo i tehdy, když je exponent v čitateli
menší nebo stejný, \textbf{definujeme} další dvě mocniny.

\begin{definiceBox}[title={Definice: nultá mocnina}]
Pro každé $a\in\mathbb R\setminus\{0\}$ definujeme
$\boxed{a^0=1}$. Výraz $0^0$ v této výuce \textbf{nedefinujeme}.
\end{definiceBox}
\begin{definiceBox}[title={Definice: záporný celočíselný exponent}]
Pro každé $a\in\mathbb R\setminus\{0\}$ a $n\in\mathbb N=\{1,2,\ldots\}$
definujeme $\boxed{a^{-n}=\dfrac1{a^n}}$.
Záporný exponent znamená \textbf{převrácenou hodnotu}, ne záporné znaménko výsledku.
\end{definiceBox}

\poznamka{Zeptej se: \uv{Proč nesmí být $a=0$?} Protože bychom dělili nulou.
Změnu $a^2:a^5$ na $a^{-3}$ nedokazujeme starou definicí, ale \emph{motivujeme}
novou definici tak, aby se zachovala známá pravidla.}

\blok{Tabulka pravidel pro mocniny: přidáváme další dva řádky}
\begin{center}\scriptsize
\begin{tabularx}{\linewidth}{@{}>{\raggedright\arraybackslash}p{33mm}>{\raggedright\arraybackslash}p{49mm}X@{}}
\toprule
\textbf{Pravidlo} & \textbf{Zápis} & \textbf{Podmínky}\\\midrule
První mocnina & $a^1=a$ & $a\in\mathbb R$\\
Nula a jednička & $0^n=0$, $1^n=1$ & $n\in\mathbb N$\\
Znaménka & $a^n>0$ pro $a>0$; pro $a<0$ kladné / záporné & $n$ sudé / liché\\
Základ $-1$ & $(-1)^n=1$ / $-1$ & $n$ sudé / liché\\
Součin stejného základu & $a^r a^s=a^{r+s}$ & $r,s\in\mathbb N$\\
Podíl stejného základu & $a^r/a^s=a^{r-s}$ & $a\ne0$, dosud $r>s\ge1$\\
Stejné exponenty & $a^r/a^r=1$ & $a\ne0$\\
Mocnina mocniny & $(a^r)^s=a^{rs}$ & $r,s\in\mathbb N$\\
Mocnina součinu & $(ab)^r=a^rb^r$ & $r\in\mathbb N$\\
Mocnina podílu & $(a/b)^r=a^r/b^r$ & $b\ne0$, $r\in\mathbb N$\\\midrule
\textbf{NOVĚ: nultá mocnina} & $\boldsymbol{a^0=1}$ & $\boldsymbol{a\ne0}$\\
\textbf{NOVĚ: záporný exponent} & $\boldsymbol{a^{-n}=1/a^n}$ & $\boldsymbol{a\ne0,\ n\in\mathbb N}$\\\bottomrule
\end{tabularx}\normalsize\end{center}
Odteď platí vzorce pro součin, podíl, mocninu mocniny, součinu a podílu
i pro \textbf{celé exponenty}, pokud jsou všechny výrazy definované.
\textbf{Obecně tedy} $a^r/a^s=a^{r-s}$ pro $a\ne0$ a $r,s\in\mathbb Z$.

\begin{prikladReseniBox}[title={Příklad 1: co udělá záporný exponent?}]
Vypočti $2^{-3}$, $(-3)^{-2}$, $-3^{-2}$,
$\left(\frac25\right)^{-2}$ a $5^0$.
\tcblower
$2^{-3}=\frac18$; $(-3)^{-2}=\frac1{(-3)^2}=\frac19$;
$-3^{-2}=-\frac19$;
$\left(\frac25\right)^{-2}=\left(\frac52\right)^2=\frac{25}{4}$;
$5^0=1$. Vždy rozlišuj \emph{základ v závorkách} a mínus před mocninou.
\end{prikladReseniBox}

\begin{zadaniReseniBox}[title={Samostatná práce 1: rychlé výpočty},breakable=false]
\begin{enumerate}[label=\alph*)]
\item $10^{-3}$; \item $(-2)^{-5}$; \item $(-4)^{-2}$;
\item $-2^{-4}$; \item $\left(-\frac12\right)^{-3}$;
\item $0^5$ a $0^0$; \item $\left(-\frac35\right)^{-2}$;
\item $(-1)^{-2027}$.
\end{enumerate}
\tcblower
 a) $\frac1{1000}$; b) $-\frac1{32}$; c) $\frac1{16}$;
 d) $-\frac1{16}$; e) $-8$; f) $0$, druhý výraz není definován;
 g) $\frac{25}{9}$; h) $-1$.
\end{zadaniReseniBox}

\blok{KISS: počítáme s celými exponenty}
\begin{prikladReseniBox}[title={Příklad 2: záporný exponent v součinu a podílu}]
Uprav $a^{-3}a^7$ a $\dfrac{x^{-2}}{x^3}$. Uveď podmínky.
\tcblower
Pro $a\ne0$: $a^{-3+7}=a^4$. Pro $x\ne0$:
$x^{-2-3}=x^{-5}=\frac1{x^5}$. Podmínky nezmizí ani po úpravě.
\end{prikladReseniBox}
\begin{zadaniReseniBox}[title={Samostatná práce 2: přepiš bez záporných exponentů}]
\begin{enumerate}[label=\alph*)]
\item $x^{-2}x^6$; \item $\frac{y^4}{y^7}$;
\item $a^{-3}b^2$; \item $\frac{p^{-2}}{p^{-5}}$;
\item $\frac{x^{-4}y^3}{x^2y^{-2}}$;
\item $\frac{m^3n^{-2}}{m^{-1}n^4}$.
\end{enumerate}
\tcblower
 a) $x^4$ ($x\ne0$); b) $1/y^3$ ($y\ne0$);
 c) $b^2/a^3$ ($a\ne0$); d) $p^3$ ($p\ne0$);
 e) $y^5/x^6$ ($x,y\ne0$); f) $m^4/n^6$ ($m,n\ne0$).
\end{zadaniReseniBox}
\begin{prikladReseniBox}[title={Příklad 3: záporný exponent celé závorky}]
Uprav $(2a^{-1}b^2)^{-2}$, předem uveď podmínky.
\tcblower
$a,b\ne0$. Použijeme mocninu součinu a mocninu mocniny:
$(2a^{-1}b^2)^{-2}=2^{-2}a^2b^{-4}=\boxed{\frac{a^2}{4b^4}}$.
\end{prikladReseniBox}
\begin{zadaniReseniBox}[title={Samostatná práce 3: mocnina mocniny},breakable=false]
\begin{enumerate}[label=\alph*)]
\item $(3^{-2})^3$; \item $(x^{-3})^{-2}$;
\item $(a^2b^{-1})^{-2}$; \item $\left(\frac{2x}{y^2}\right)^{-2}$;
\item $(2x)^0$ a $2x^0$ pro $x\ne0$.
\end{enumerate}
\tcblower
 a) $3^{-6}=1/729$; b) $x^6$ ($x\ne0$);
 c) $b^2/a^4$ ($a,b\ne0$);
 d) $y^4/(4x^2)$ ($x,y\ne0$);
 e) $1$ a $2$, obojí při $x\ne0$.
\end{zadaniReseniBox}
\poznamka{Pozor: $x^{-3}$ není záporné číslo automaticky; záleží na základu.
V posledním příkladu rozliš celý základ $(2x)^0$ a samostatnou mocninu $x^0$.
Pro $x=0$ je $2x^0$ nedefinováno, nikoli $2$.}

\blok{Rozšíření a rezerva}
\begin{zadaniReseniBox}[title={Rozšiřující příklad 1: znaménka a podmínky}]
\begin{enumerate}[label=\alph*)]
\item Uprav $\dfrac{(2a^{-2}b)^3}{4a^{-1}b^{-2}}$.
\item Jakou hodnotu má $(x-1)^0$ a kdy neexistuje?
\item Uprav $\dfrac{2^{3n+1}4^{2n-2}}{8^{n-1}}$ pro $n\in\mathbb Z$.
\end{enumerate}
\tcblower
 a) $2b^5/a^5$ pro $a,b\ne0$; b) $1$ pro $x\ne1$,
 pro $x=1$ vznikne $0^0$; c) $2^{4n}$, neboť všechny základy jsou mocniny dvojky.
\end{zadaniReseniBox}
\begin{zadaniReseniBox}[title={Rozšiřující příklad 2: najdi a oprav chybu}]
\begin{enumerate}[label=\alph*)]
\item Žák píše $(-2)^{-4}=-\frac1{16}$. Proč je to špatně?
\item Žák píše $(ab)^{-2}=a^{-2}+b^{-2}$. Oprav.
\item Je $\dfrac{x^2}{x^2}=1$ i pro $x=0$?
\end{enumerate}
\tcblower
 a) Sudá mocnina záporného základu je kladná: $1/16$.
 b) Správně $a^{-2}b^{-2}=1/(a^2b^2)$ při $a,b\ne0$.
 c) Ne, původní výraz není pro $x=0$ definován.
\end{zadaniReseniBox}

\blok{Domácí procvičení}
\begin{zadaniReseniBox}[title={Domácí úkol 1: záporné exponenty}]
\begin{enumerate}[label=\alph*)]
\item Vypočti $(-3)^{-3}$ a $\left(\frac25\right)^{-2}$.
\item Uprav $x^{-4}x^7$ a $\dfrac{y^2}{y^{-3}}$.
\item Uprav $(3a^{-2}b)^ {-2}$.
\item Uprav $\dfrac{6x^{-2}y^3}{3x^4y^{-2}}$.
\item Urči podmínky a hodnotu $(t+2)^0$.
\end{enumerate}
\tcblower
 a) $-1/27$, $25/4$; b) $x^3$ ($x\ne0$), $y^5$ ($y\ne0$);
 c) $a^4/(9b^2)$ ($a,b\ne0$);
 d) $2y^5/x^6$ ($x,y\ne0$);
 e) $1$ pro $t\ne-2$, jinak nedefinováno.
\end{zadaniReseniBox}
\end{document}
